\begin{theorem}[Horizontal sections of sampled polygons]
	\label{thm:area_law_sample_rows}
	\lean{
		TNLean.PEPS.AreaLaw.Geometry.latticeRow,
		TNLean.PEPS.AreaLaw.Geometry.mem_latticeRow,
		TNLean.PEPS.AreaLaw.Geometry.TemplatePolygon.convex_region,
		TNLean.PEPS.AreaLaw.Geometry.TemplatePolygon.isCompact_region,
		TNLean.PEPS.AreaLaw.Geometry.horizontal_mem_of_mem_of_le,
		TNLean.PEPS.AreaLaw.Geometry.TemplatePolygon.isCompact_horizontalSection,
		TNLean.PEPS.AreaLaw.Geometry.TemplatePolygon.horizontalSection_eq_Icc,
		TNLean.PEPS.AreaLaw.Geometry.Template.latticeRow_eq_Icc_ceil_floor,
		TNLean.PEPS.AreaLaw.Geometry.Template.mem_sample_of_row_between,
		TNLean.PEPS.AreaLaw.Geometry.Template.latticeRow_eq_Icc,
		TNLean.PEPS.AreaLaw.Geometry.ceil_affine_row,
		TNLean.PEPS.AreaLaw.Geometry.floor_affine_row
	}
	\leanok
	\uses{def:area_law_template}
	Let $P$ be one of the real polygons of a template, put
	$S=P\cap\mathbb Z^2$, and for $y\in\mathbb Z$ write
	$R_y(S)=\{x\in\mathbb Z:(x,y)\in S\}$.
	Every nonempty real horizontal section of $P$ is a closed interval
	$[a,b]$ whose lattice intersection rounds the endpoints:
	\begin{align}
		R_y(S) & =[\lceil a\rceil,\lfloor b\rfloor]\cap\mathbb Z.
	\end{align}
	When this integer interval is nonempty, its endpoints are its attained
	minimum and maximum; for every real $c$ and integers $m,y$, rounding
	commutes with integer shifts:
	\begin{align}
		\lceil c+my\rceil   & =\lceil c\rceil+my,   \\
		\lfloor c+my\rfloor & =\lfloor c\rfloor+my.
	\end{align}
\end{theorem}
\begin{proof}
	\leanok
	Because a finite convex hull is compact and convex, its intersection with
	a horizontal line is a compact interval connecting its attained extrema.
	Intersecting that interval with $\mathbb Z$ rounds the lower endpoint
	upward and the upper endpoint downward, and integer translation commutes
	with both floor and ceiling, establishing the initial row step in the
	proof of Lemma~9.4 of \cite{OpenAI2026AreaLaw}.
\end{proof}

\begin{theorem}[Area of a template]
	\label{thm:area_law_template_area}
	\lean{
		TNLean.PEPS.AreaLaw.Geometry.Template.sample_subset_box,
		TNLean.PEPS.AreaLaw.Geometry.Template.card_sample_le,
		TNLean.PEPS.AreaLaw.Geometry.template_card_le
	}
	\leanok
	\uses{def:area_law_template}
	If $C_{\rm tpl}\ge1$, a template satisfies $|T|\le9ns_0$.
\end{theorem}
\begin{proof}
	\leanok
	Each nonempty sample lies in an integer square of radius $s_0$ centered at
	any one of its sites and therefore contains at most $(2s_0+1)^2$ sites,
	whereas empty samples contribute zero. Summing over the $N_0$ pieces and
	using $(2s_0+1)^2\le9s_0(s_0+1)$ together with $N_0(s_0+1)\le n$ yields
	$|T|\le9ns_0$.
\end{proof}

\begin{theorem}[Rows and unions under ambient dilation]
	\label{thm:area_law_dilation_rows_union}
	\lean{
		TNLean.PEPS.AreaLaw.mem_ambientDilation_iff,
		TNLean.PEPS.AreaLaw.ambientDilation_mono,
		TNLean.PEPS.AreaLaw.ambientDilation_mono_radius,
		TNLean.PEPS.AreaLaw.ambientDilation_add,
		TNLean.PEPS.AreaLaw.ambientDilation_biUnion_sdiff_subset,
		TNLean.PEPS.AreaLaw.card_ambientDilation_biUnion_sdiff_le,
		TNLean.PEPS.AreaLaw.Geometry.mem_latticeRow_ambientDilation,
		TNLean.PEPS.AreaLaw.Geometry.Template.mem_latticeRow_dilation_iff
	}
	\leanok
	\uses{def:area_law_template,thm:area_law_sample_rows}
	For a finite $S\subset\mathbb Z^2$, let
	$S_r=S+([-r,r]^2\cap\mathbb Z^2)$, so a point belongs to $S_r$ if and only
	if both of its coordinate differences from some point of $S$ are at most
	$r$ in absolute value. Ambient dilation is monotone in both $S$ and $r$,
	satisfies $(S_r)_t=S_{r+t}$, and has horizontal rows
	\begin{align}
		R_y(S_r) & =
		\bigcup_{z=y-r}^{y+r}\ \bigcup_{a\in R_z(S)}
		                       ([a-r,a+r]\cap\mathbb Z),
	\end{align}
	where each nonempty row $R_z(S)$ of a sampled polygon contributes the
	interval from its minimum minus $r$ to its maximum plus $r$.
	When $S=\bigcup_{i\in I}S_i$ is a finite union of possibly empty,
	overlapping, or disconnected pieces, every $r,t\ge0$ satisfy
	\begin{align}
		S_r\setminus S_t   & \subseteq\bigcup_{i\in I}((S_i)_r\setminus(S_i)_t), \\
		|S_r\setminus S_t| & \le\sum_{i\in I}|(S_i)_r\setminus(S_i)_t|.
	\end{align}
\end{theorem}
\begin{proof}
	\leanok
	The coordinate bounds defining $S_r$ decompose directly into vertical and
	horizontal intervals to give the row formula. For a finite union
	$S=\bigcup_{i\in I}S_i$, any point in $S_r\setminus S_t$ lies within
	distance $r$ of some constituent $S_i$ while avoiding $(S_i)_t$, which
	gives the set inclusion and, by union subadditivity, the cardinality bound.
\end{proof}

\begin{theorem}[Supporting strips of allowed polygons]
	\label{thm:area_law_polygon_strips}
	\leanok
	\lean{
		TNLean.PEPS.AreaLaw.Geometry.TemplatePolygon.vertices,
		TNLean.PEPS.AreaLaw.Geometry.TemplatePolygon.region_eq_convexHull_vertices,
		TNLean.PEPS.AreaLaw.Geometry.TemplatePolygon.mem_region_iff_normal_bounds,
		TNLean.PEPS.AreaLaw.Geometry.TemplatePolygon.exists_four_strip_bounds
	}
	\uses{def:area_law_template}
	Each allowed polygon equals the intersection of the supporting strips for
	the four linear forms $x,y,x+y,x-y$, or equivalently the intersection of all
	closed supporting slabs with allowed normal directions.
\end{theorem}
\begin{proof}
	\leanok
	Every supporting slab contains the convex hull of the vertices, and the
	extrema of each linear form over the finite vertex set are attained.
	Conversely, for a triangle, dividing the three oriented side functionals by
	its nonzero signed determinant (of either sign) produces nonnegative
	barycentric coordinates, while for a rectangle the orthogonal side
	coordinates in $[0,1]$ produce four nonnegative convex weights; since every
	allowed normal is a scalar multiple of one of $x,y,x+y,x-y$, the four
	supporting strips cut out the polygon.
\end{proof}

% Auxiliary to Definition 9.3 (scanner:template), 08-scanner.tex, lines 546--569.
% Revision: adc7f1241b42e322a6451854ab7e4b4c146bf78a.
\begin{theorem}[Finite integer points in allowed polygons]
	\label{thm:area_law_polygon_integer_points_finite}
	\lean{TNLean.PEPS.AreaLaw.Geometry.TemplatePolygon.finite_integerPoints}
	\leanok
	\uses{def:area_law_template}
	Let $P\subseteq\mathbb R^2$ be a closed nondegenerate triangle or
	rectangle whose sides have slopes in $\{0,\infty,1,-1\}$.
	Then $P\cap\mathbb Z^2$ is finite.
\end{theorem}

\begin{proof}
	\leanok
	\uses{thm:area_law_polygon_strips}
	By Theorem~\ref{thm:area_law_polygon_strips}, there exist
	$l_x,u_x,l_y,u_y\in\mathbb R$ such that
	$P\subseteq[l_x,u_x]\times[l_y,u_y]$, which implies
	\begin{align}
		P\cap\mathbb Z^2 & \subseteq
		                   ([\lceil l_x\rceil,\lfloor u_x\rfloor]\cap\mathbb Z)
		\times([\lceil l_y\rceil,\lfloor u_y\rfloor]\cap\mathbb Z).\notag
	\end{align}
	Because each integer interval is finite, their Cartesian product and its
	subset $P\cap\mathbb Z^2$ are finite.
\end{proof}

% Auxiliary to Definition 9.3 (scanner:template) and 10-geometry.tex,
% lines 299--318 and 732--740.
% Revision: adc7f1241b42e322a6451854ab7e4b4c146bf78a.
\begin{theorem}[A template for the integer points of a fan run]
	\label{thm:area_law_fan_run_template}
	\lean{TNLean.PEPS.AreaLaw.Geometry.CellFanRun.exists_template}
	\leanok
	\uses{def:area_law_template,def:area_law_dyadic_cells,def:area_law_cell_fan,
		def:area_law_fan_runs}
	Let $C_{\mathrm{tpl}},C_B\in\mathbb R$ satisfy
	$C_{\mathrm{tpl}}\ge1$ and $C_B\ge16C_{\mathrm{tpl}}$.
	Fix any origin $o\in\mathbb R^2$, scale $\ell\in\mathbb N$, and
	cell index $z\in\mathbb Z^2$, together with midpoint subdivisions and a
	two-coloring of the fan triangles of the closed dyadic square of side
	$u=2^\ell$. For any run $R$ whose lattice intersection
	$T_R\cap\mathbb Z^2$ is nonempty, there exists a template $T$ with
	parameters $s_0=2^\ell$ and $n=\lceil C_Bu\rceil$, formed by at most eight
	polygonal pieces, such that $x\in T$ if and only if $x\in T_R$ for every
	$x\in\mathbb Z^2$.
\end{theorem}

\begin{proof}
	\leanok
	\uses{def:area_law_template,def:area_law_fan_runs,
		thm:area_law_cell_fan_count,thm:area_law_cell_fan_cover,
		thm:area_law_fan_run_cover,
		thm:area_law_polygon_integer_points_finite,
		thm:area_law_dyadic_cell_distance}
	Select the $m$ fan triangles $P_e$ indexed by $R$. Because
	$T_R\cap\mathbb Z^2$ is nonempty, at least one such triangle contains a
	lattice point, while Theorem~\ref{thm:area_law_cell_fan_count} gives
	$1\le m\le8$. Assigning to each piece $P_e$ its finite lattice sample
	$P_e\cap\mathbb Z^2$ via
	Theorem~\ref{thm:area_law_polygon_integer_points_finite} (allowing empty
	individual samples) yields a union equal to $T_R\cap\mathbb Z^2$ across all
	shared edges.

	By Theorem~\ref{thm:area_law_cell_fan_cover}, every piece lies in the same
	closed dyadic square of side $u$, so
	Theorem~\ref{thm:area_law_dyadic_cell_distance} with equal cell indices and
	radius zero bounds the maximum-norm diameter of each piece and of the total
	sample by $u=s_0$. Since $u\ge1$ and $C_B\ge16$, both $s_0$ and $n$ are
	positive and satisfy
	\begin{align}
		C_{\mathrm{tpl}}m(s_0+1) & \le8C_{\mathrm{tpl}}(u+1)
		\le16C_{\mathrm{tpl}}u\le C_Bu\le n,\notag
	\end{align}
	verifying all scale and diameter conditions of a template.
\end{proof}

\begin{theorem}[Rounded affine row profiles]
	\label{thm:area_law_affine_row_profiles}
	\leanok
	\lean{
		TNLean.PEPS.AreaLaw.Geometry.Template.exists_sample_four_strip_bounds,
		TNLean.PEPS.AreaLaw.Geometry.Template.exists_latticeRow_profile,
		TNLean.PEPS.AreaLaw.Geometry.Template.latticeRow_nonempty_between,
		TNLean.PEPS.AreaLaw.Geometry.Template.latticeRow_min_step,
		TNLean.PEPS.AreaLaw.Geometry.Template.latticeRow_max_step,
		TNLean.PEPS.AreaLaw.Geometry.Template.exists_sample_in_window_between,
		TNLean.PEPS.AreaLaw.Geometry.Template.exists_nearby_sample_in_row
	}
	\uses{thm:area_law_polygon_strips,thm:area_law_sample_rows}
	For each sampled polygon $S$ there exist eight integers
	$l_x,u_x,l_y,u_y,l_+,u_+,l_-,u_-$ such that $R_y(S)$ is empty outside
	$[l_y,u_y]$ and, for $y\in[l_y,u_y]$,
	\begin{align}
		R_y(S) & =
		[\max(l_x,l_+-y,l_-+y),\min(u_x,u_+-y,u_-+y)]\cap\mathbb Z.
	\end{align}
	The nonempty rows $R_y(S)$ form an integer interval of $y$-values across
	which the attained lower and upper endpoints each change by at most one per
	row step; restricting to a vertical window preserves consecutive occupied
	columns and allows any sampled point to shift to an occupied adjacent row
	with horizontal displacement at most one.
\end{theorem}
\begin{proof}
	\leanok
	Rounding the real lower strip bounds upward and the upper bounds downward
	commutes with integer shifts in $y$, expressing the lower row endpoint as
	the maximum of three affine functions of slopes $0,-1,1$ and the upper
	endpoint as the minimum of three such functions. Pairwise comparisons
	between these lower and upper affine branches carve out an interval of
	nonempty rows, and because each affine branch changes by at most one per
	unit step in $y$, so do their maximum and minimum.
\end{proof}

\begin{theorem}[Single-piece ambient dilation layers]
	\label{thm:area_law_piece_depth_layers}
	\leanok
	\lean{
		TNLean.PEPS.AreaLaw.Geometry.Template.mem_latticeRow_dilation_of_between,
		TNLean.PEPS.AreaLaw.Geometry.Template.latticeRow_dilation_eq_Icc,
		TNLean.PEPS.AreaLaw.Geometry.Template.latticeRow_dilation_succ_bounds,
		TNLean.PEPS.AreaLaw.Geometry.Template.card_latticeRow_layer_le_four,
		TNLean.PEPS.AreaLaw.Geometry.Template.card_latticeRow_dilation_le,
		TNLean.PEPS.AreaLaw.Geometry.Template.exists_dilation_row_domain,
		TNLean.PEPS.AreaLaw.Geometry.Template.card_dilation_layer_le
	}
	\uses{thm:area_law_affine_row_profiles,thm:area_law_dilation_rows_union,
		thm:area_law_template_area}
	For any sampled piece $S$, each dilated horizontal row $R_y(S_r)$ is an
	interval of at most $2s_0+2r+1$ sites, and increasing the radius from $r$
	to $r+1$ adds at most four sites to each nonempty row of $S_r$. When $S$
	is nonempty, there exist integers $a\le b$ with $b-a\le2s_0$ such that the
	occupied rows of $S_r$ are $[a-r,b+r]\cap\mathbb Z$, and for every
	$1\le j\le s_0$ (including when $S$ is empty),
	\begin{align}
		|S_j\setminus S_{j-1}| & \le24(s_0+1).
	\end{align}
\end{theorem}
\begin{proof}
	\leanok
	Intersecting $S$ with a vertical radius-$r$ window yields consecutive
	occupied columns whose horizontal $r$-neighborhood is a single interval of
	width bounded by $2s_0+2r+1$. Passing from radius $j-1$ to radius $j$
	shifts any clamped row index by at most one vertical step (moving the
	horizontal endpoints by at most one) and widens the horizontal radius by one,
	so each previously occupied row gains at most two sites per end (four in
	total). Summing at most four new sites across the old occupied rows and
	adding the two newly occupied extreme rows bounds $|S_j\setminus S_{j-1}|$
	by $12s_0+12j-2\le24(s_0+1)$.
\end{proof}

\begin{theorem}[Template depth layers with an explicit constant]
	\label{thm:area_law_template_depth_layers}
	\leanok
	\lean{TNLean.PEPS.AreaLaw.Geometry.template_layer_card_le}
	\uses{def:area_law_template,thm:area_law_piece_depth_layers,
		thm:area_law_dilation_rows_union}
	If $C_{\rm tpl}\ge24$ and $1\le j\le s_0$, the template satisfies
	\begin{align}
		|T_j\setminus T_{j-1}| & \le n.
	\end{align}
\end{theorem}
\begin{proof}
	\leanok
	Summing the single-piece layer bound across the $N_0$ pieces (whether
	overlapping, empty, or disconnected) and applying the template scale bound
	gives
	\begin{align}
		|T_j\setminus T_{j-1}|
		 & \le24N_0(s_0+1)
		\le C_{\rm tpl}N_0(s_0+1)\le n.
	\end{align}
\end{proof}

These geometric estimates govern only the lattice cardinalities of templates
and their dilations; the entropy and physical area-law bounds follow in the
subsequent sections.
